\documentclass[twocolumn]{iopjournal}

\usepackage{amsmath,amssymb,amsfonts}

\fancyhfoffset[L]{0mm}
\setlength{\marginparwidth}{0pt}
\setlength{\marginparsep}{0pt}

% ==========================================
% ZENODO PREPRINT HEADERS (ACTIVE)
% ==========================================
\renewcommand{\articletype}[1]{%
  {\vspace*{-8mm}\noindent \Large \sf Preprint}\par
  \vspace*{6mm}\noindent{\scriptsize \sf{\bfseries \MakeUppercase{#1}}}}

\fancyhead[L]{}
\fancyhead[R]{}

\begin{document}

\articletype{Supplementary Material}

\title{Supplementary Material: Chiral Holonomy\\
of the Closed Tetrad Frame}

\author{Stephen Nelson}

\affil{Field Dynamics Research}

\ead{stephen@fielddynamics.org}

\begin{abstract}
The companion letter~\cite{Nelson2026b} proved that the
24~torsion components of the $K = 4$ tetrad connection
admit the factorisation $24 = D \times 2K = 3 \times 8$,
exclusively at $K = 4$.  This note identifies the two
factors with the geometry of the closed frame: $2K = 8$
is the count of null directions~\cite{Nelson2026a}
(vertices of the stellated octahedron) and $D = 3$ is
the count of transverse rotation planes per direction.  The parity
distinction between TetA and TetB, established
in~\cite{Nelson2026b} at the partition level, resolves at
the component level into opposite chirality: the
dualized torsion vector reverses sign between the two
sub-frames.  A direct computation of parallel transport
around all twelve Hamiltonian cycles of the cube graph
shows that uniform transport (ignoring chirality)
produces trivial holonomy, while chiral transport
(respecting the TetA/TetB sign distinction) produces
nontrivial holonomy in every case.  The chirality
carried by the torsion sector is therefore the mechanism
that generates rotational closure of the gravitational
field.
\end{abstract}

% ==================================================================

\section*{The two factors}

The companion letter~\cite{Nelson2026b} proved
(Section~5, Theorem) that the torsion count
$K \times K(K{-}1)/2 = 24$ admits the factorisation
$D \times 2K$ exclusively at $K = 4$, where $D = K - 1$.
The algebraic identity was established there.  This
section identifies what the two factors refer to in the
geometry of the closed frame.

The stellated octahedron~\cite{Nelson2026a} has
$2K = 8$ vertices: four from TetA (the identity
component) and four from TetB (the parity component).
In Cartesian coordinates, these are
$(\pm 1/\sqrt{3},\, \pm 1/\sqrt{3},\, \pm 1/\sqrt{3})$
--- the eight vertices of the cube dual to the
octahedral core.  Each vertex lies on the celestial
sphere $S^2$ and defines a null direction of the
frame~\cite{Nelson2026a}.
The factor $2K$ in the torsion factorisation is the
count of null directions of the closed frame.

The cube has four body diagonals (antipodal pairs).
Each vertex belongs to exactly one diagonal and is
therefore transverse to the remaining three.  Each
transverse body diagonal defines a rotation plane
perpendicular to it; three transverse diagonals
therefore give three independent rotation planes, one
per spatial dimension.  The factor $D$ in
the torsion factorisation is the count of transverse
rotation planes per null direction.

The $24$~torsion components therefore correspond to
the set of (null direction, rotation plane) pairs:
\begin{equation}
  24 \;=\; \underbrace{8}_{\text{null directions}}
        \;\times\; \underbrace{3}_{\text{rotation planes}}
  \;=\; 2K \times D.
  \label{eq:24}
\end{equation}

% ==================================================================

\section*{The chirality split}

The companion letter~\cite{Nelson2026b} proved that
TetA is the identity component and TetB is the parity
component under $P: e^i \to -e^i$.  The
Lemma and Proposition of that letter establish that the
torsion sector is the necessary carrier of the TetA/TetB
distinction: a torsion-free connection, built from the
parity-invariant metric alone, cannot resolve the two
roots.

At the level of individual torsion components, the
parity distinction manifests as opposite chirality.
In the closed-frame geometry, the octahedral core has
eight faces, each carrying an attached tetrahedron
whose apex is one of the eight cube vertices (null
directions).  Each attached tetrahedron has four faces:
one base face (shared with the octahedron, the
propagation face) and three lateral faces.  The three
lateral faces at each null direction are the three
transverse rotation planes identified in the previous
section.  For each null direction, the three torsion
components can be collected into a dualized spatial
rotation vector using the Levi-Civita symbol:
\begin{equation}
  \tau^I{}_a = \tfrac{1}{2}\,
    \varepsilon_{abc}\, T^I{}_{bc}.
  \label{eq:dual}
\end{equation}
This $\tau^I{}_a$ is a spatial 3-vector: one component
per lateral face, one per transverse rotation plane.

Under the orientation reversal that exchanges TetA and
TetB, the Levi-Civita symbol reverses sign:
$\varepsilon^{\text{B}}_{abc} =
-\varepsilon^{\text{A}}_{abc}$, because the spatial
triad reverses handedness.  Therefore:
\begin{equation}
  \tau^I{}_a\big|_{\text{TetB}}
  = -\,\tau^I{}_a\big|_{\text{TetA}}.
  \label{eq:flip}
\end{equation}
The four TetA null directions carry rotation vectors
with chirality~$(+)$.  The four TetB null directions
carry rotation vectors with chirality~$(-)$.  The
24~components split as $12^{+} + 12^{-}$.

% ==================================================================

\section*{The holonomy computation}

The eight null directions form a bipartite graph: every
cube edge connects a TetA vertex to a TetB vertex
(flipping one coordinate bit changes the popcount
parity).  Each vertex has three neighbours, all in the
opposite sub-frame.  The angular separation across each
edge is $\arccos(1/3) \approx 70.53^{\circ}$.

A closed transport path through all eight directions
follows a Hamiltonian cycle on this graph, alternating
sub-frames at every step: $\text{A} \to \text{B} \to
\text{A} \to \text{B} \to \cdots$ (eight steps).
At each step from vertex~$v_i$ to vertex~$v_j$, the
frame rotates by angle
$\theta = \arccos(v_i \cdot v_j)$ about the axis
$\hat{n} = (v_i \times v_j)/|v_i \times v_j|$.  The
chirality of the source vertex determines the sign:
$+\theta$ from TetA, $-\theta$ from TetB.

This sign assignment is not a convention; it is a
consequence of the chirality split established in the
previous section.  The torsion at a TetA vertex
carries chirality~$(+)$; the torsion at a TetB vertex
carries chirality~$(-)$.  Unlike curvature, which
contributes to holonomy at second order in the loop
area, torsion contributes at first order in the path
length~\cite{Hehl1976}.  The rotation accumulated at
each step therefore inherits the sign of the source
vertex's torsion directly.

\medskip
\noindent\textbf{Proposition.}
\textit{Let $H_{\mathrm{uniform}}$ denote the holonomy
of a Hamiltonian cycle on the cube graph with rotation
angle $+\theta$ at every step, and let
$H_{\mathrm{chiral}}$ denote the holonomy of the same
cycle with the sign of $\theta$ determined by the source
chirality ($+$ for TetA, $-$ for TetB).  Then:}

\begin{itemize}
\item[(a)] \textit{$H_{\mathrm{uniform}} = \mathbf{1}$
for all twelve Hamiltonian cycles of the cube.}
\item[(b)] \textit{$H_{\mathrm{chiral}} \neq \mathbf{1}$
for all twelve Hamiltonian cycles of the cube.}
\end{itemize}

\medskip\noindent\textit{Proof.}
The cube graph on eight vertices admits exactly twelve
Hamiltonian cycles (up to starting vertex).  Both claims
are verified by direct computation: at each step, the
$3 \times 3$ rotation matrix is constructed via
Rodrigues' formula,
\begin{equation}
  R(\hat{n}, \alpha) = I + \sin\alpha\,[\hat{n}]_{\times}
  + (1 - \cos\alpha)\,[\hat{n}]_{\times}^2,
  \label{eq:rodrigues}
\end{equation}
where $\alpha = \pm\theta$ and $[\hat{n}]_{\times}$ is
the skew-symmetric matrix of $\hat{n}$.  The eight
matrices are multiplied in path order.

For part~(a): all twelve products yield the
$3 \times 3$ identity matrix to machine precision.  The
exact cancellation is a consequence of the octahedral
symmetry of the cube: uniform rotation at each step
respects the full symmetry group.

For part~(b): the sign alternation breaks the octahedral
symmetry.  The positive and negative rotations act about
\emph{different} axes (each axis $v_i \times v_j$
depends on both endpoints), so reversing the sign at
alternating steps does not invert those steps but
produces a qualitatively different composition.

All twelve chiral products yield nontrivial rotations:
six cycles produce a net rotation of $135^{\circ}$
$(= 3\pi/4)$ and six produce $155.8^{\circ}$
$(\approx 7\pi/8)$.  The two classes correspond to the
two orientation classes of Hamiltonian cycles on the
cube.

As an explicit example, the cycle
$v_1 \!\to\! v_6 \!\to\! v_3 \!\to\! v_5 \!\to\! v_4
\!\to\! v_7 \!\to\! v_2 \!\to\! v_8 \!\to\! v_1$
(with vertex labelling as in~\cite{Nelson2026a}) gives
$H_{\mathrm{uniform}} = \mathbf{1}$ and
$H_{\mathrm{chiral}} = R(\hat{m},\, 135^{\circ})$, a
rotation by $135^{\circ}$ about the axis $\hat{m}$.
\hfill$\square$

\medskip
A self-contained verification script (approximately
forty lines of Python using NumPy) is provided as
supplementary code.

% ==================================================================

\section*{The vortex}

The holonomy per cycle is a rotation by a fixed angle
$\Theta \neq 0$ about a fixed axis.  In the
closed-frame geometry, the eight null directions (cube
vertices) are connected through the octahedral core:
every transition from one null direction to the next
crosses a face of the core octahedron.  The field
origin~$\mathcal{O}$, as the centroid of the
octahedron, lies on every such crossing path.  The
transport therefore passes through~$\mathcal{O}$ at
every $\text{TetA} \leftrightarrow \text{TetB}$
transition.  A path that advances through a common
point while rotating by a constant angle per cycle is
helical.

The rotation axis defines a preferred spatial direction.
A nontrivial element of $\mathrm{SO}(3)$ has a
one-dimensional fixed subspace (the rotation axis) and a
unique two-dimensional invariant subspace (the plane
perpendicular to it); these are the only invariant
subspaces.  Transport organised by such a
rotation therefore accumulates in the perpendicular plane.
Conversely, the uniform case (Proposition, part~1) gives
identity holonomy for all twelve cycles: no axis is
selected, and spherical symmetry is preserved.  The
chiral case breaks that symmetry by selecting exactly one
rotation plane.  The closure geometry is
planar---perpendicular to the holonomy axis---not
spherically symmetric.  The holonomy angles computed
above ($135^{\circ}$, $155.8^{\circ}$) represent the
torsion sector at full strength: the asymptotic regime
where $1 - \mu \approx 1$.  The constitutive law
modulates this strength radially.  At the galactic
centre, $x \to \infty$ and $1 - \mu \to 0$; the torsion
fraction is asymptotically suppressed and the Christoffel
sector dominates.  The chiral holonomy is therefore
active at the rim and suppressed at the centre.  The
result is a spherically symmetric Christoffel-dominated
throat transitioning smoothly into a planar
torsion-dominated disk---exactly the bulge-to-disk
structure rendered in the main letter's Figure~1.  The
global coherence of the disk axis follows from the
smoothness of the connection: on a non-singular manifold
the affine connection parallel-transports frames
continuously from point to point, so where the torsion
sector is active its holonomy axis varies smoothly rather
than randomising into disconnected domains.

This planarity is a qualitative prediction that
distinguishes the present framework from the cold dark
matter interpretation.  The observed gravitational
discrepancy at galactic scales is not in dispute: both
frameworks reproduce the radial acceleration relation
and the function $\mu(x) = x/(x{+}1)$.  They differ on
what the discrepancy represents.  The cold dark matter
framework interprets it as an invisible mass component,
typically modelled with approximately spherical
halos.  The closed tetrad frame identifies it
as the field response of the second tetrahedral
chamber (TetB), carried by the 24~torsion components.
This response is not mass; it is topology.  Its spatial
distribution is not fitted but determined: the chiral
holonomy defines an axis, and the resulting geometry is
a disk.  Observed spiral galaxies are disks.

The constitutive law $\mu(x) = x/(x{+}1)$, derived from
the exchange symmetry in~\cite{Nelson2026b}, is the
radial profile of the vortex: TetA-dominated
($\mu \to 1$) at the inner throat, TetB-dominated
($\mu \to 0$) at the outer rim, with the transition at
$x = 1$.  The constitutive law determines the algebraic
partition.  The chiral holonomy determines the spatial
geometry.  Together they define the complete
gravitational field.

% ==================================================================

\section*{Summary}

\vspace{6pt}
\noindent
{\footnotesize
\begin{tabular*}{\columnwidth}{@{\extracolsep{\fill}}lll@{}}
\hline
 & Letter [2] & This note \\
\hline
$24 = D \times 2K$ & algebraic identity & geometric identification \\
$2K = 8$ & coupling range & null directions \\
$D = 3$ & spatial dimension & transverse planes \\
Chirality & parity of partition & sign of $\tau^I{}_a$ \\
Holonomy & (not addressed) & computed, nontrivial \\
Geometry & (not addressed) & vortex with axis \\
\hline
\end{tabular*}}
\vspace{6pt}

The 24~torsion components that~\cite{Nelson2026b}
identified as the carrier of TetB organise into
eight null directions, each carrying three transverse
rotation planes.  Their parity-resolved chirality
breaks the exact holonomy cancellation of the
directional scaffold, forcing the field to close as a
vortex.  The constitutive law $\mu(x) = x/(x{+}1)$,
already derived in~\cite{Nelson2026b}, is the radial
profile of this vortex.

% ==================================================================

\begin{thebibliography}{99}

\bibitem{Nelson2026a}
S.~Nelson,
\textit{The Coordinate-Space Geometry of the Tetrad
Frame},
Class.\ Quantum Grav.\ (2026), submitted (CQG-115428).

\bibitem{Nelson2026b}
S.~Nelson,
\textit{The Other Half of Gravity: Torsion and the
Two-Chamber Structure of the Tetrad Frame},
Class.\ Quantum Grav.\ (2026), submitted (CQG-115447).

\bibitem{Hehl1976}
F.~W.~Hehl, P.~von~der~Heyde, G.~D.~Kerlick and
J.~M.~Nester,
\textit{General relativity with spin and torsion:
Foundations and prospects},
Rev.\ Mod.\ Phys.\ \textbf{48} (1976) 393.

\end{thebibliography}

\end{document}
